\documentclass[11pt,a4paper]{article}

\usepackage[margin=1in]{geometry}
\usepackage{amsmath,amssymb,amsthm}
\usepackage{graphicx}
\usepackage{booktabs}
\usepackage{microtype}
\usepackage{siunitx}
\usepackage{xcolor}
\usepackage[hidelinks,colorlinks=true,linkcolor=blue!50!black,citecolor=blue!40!black,urlcolor=blue!50!black]{hyperref}
\usepackage[numbers,sort&compress]{natbib}

\sisetup{detect-all=true}

\title{\textbf{Cortical Closure Residual as a First-Order Proxy for Chronic Functional Dyscoherence}\\[2pt]
\large Synthetic Pre-Registered Validation, Trait-Anxiety Real-Data Anchor, and a Pre-Registered PTSD Replication\\[1pt]
\small (Solution Lab experiment 005, open research note --- pre-peer-review draft, metric version CCR-v1)}

\author{%
  Lee Smart\\[4pt]
  \textit{Institute of Vibrational Field Dynamics}\\[2pt]
  \texttt{contact@vibrationalfielddynamics.org}\\[2pt]
  \texttt{@vfd\_org}%
}
\date{May 2026}

\begin{document}
\maketitle
\thispagestyle{empty}

\begin{abstract}
\noindent
\emph{Life as Closure} \citep{SmartLifeAsClosure} \S 8 treats sustained dyscoherence as a P-A-conditional structural proxy for suffering: persistent distance between a living frame's substrate state and its closure-stable self-model subspace $\Sigma_\mathcal{I}$. We supply a measurable first-order empirical proxy for chronic functional dyscoherence on cortical EEG: the per-subject calibrated cortical closure residual (CCR) of \citet{SolutionLab002}, applied between-cohort to clinical-state contrasts. This is not the operator $\mathcal{D}_\mathcal{I}(v) = \|v - P_{\Sigma_\mathcal{I}}(v)\|$ itself --- $P_{\Sigma_\mathcal{I}}$ is not directly measured on any substrate in this paper. CCR is the L\textsuperscript{2}-distance from a within-cohort reference centroid in observable feature space, with the centroid acting as a cohort-aggregate proxy for $P_{\Sigma_\mathcal{I}}$. The methodology is conditional on the $D_1$--$D_5$ applicability diagnostic of the companion methodology paper \citep{SmartClosureDistance}. We report four results. (i)~\textbf{Synthetic pre-registered validation}: five preregistered hypotheses on synthetic PTSD-vs-healthy data all pass at threshold ($5/5$). (ii)~\textbf{Real-data anchor}: application to the OpenNeuro trait-anxiety cohort \texttt{ds007609} \citep{ShalamberidzeAnxiety2025} ($n=24$, low-STAI vs high-STAI quartile contrast) returns between-cohort Cohen's $d = +0.79$ (Welch $t = +1.93$, $p = 0.07$) using the SL-002 \texttt{compute\_features} pipeline verbatim with default weights. The cross-substrate applicability diagnostic predicts PASS \emph{a priori} for this substrate. Diagnostic prediction and empirical outcome agree. (iii)~\textbf{Independent-cohort replication on mood manipulation \texttt{ds004315}} \citep{CavanaghMoodManipulation}: the same pipeline applied without parameter retuning to a different paradigm (within-subject sad-mood induction vs neutral, $n = 50$, $25$ per arm) returns Cohen's $d = +0.70$ (Welch $t = +2.47$, $\mathbf{p = 0.018}$). The unpaired-design diagnostic predicted FAIL on $D_2$/$D_3$ at this cohort size; the empirical outcome is PASS at $d > 0.5$ threshold --- a diagnostic false negative consistent with the per-design threshold-calibration issue noted in \citep{SmartClosureDistance}. (iv)~\textbf{Scope refinement from companion negative}: the same methodology applied to dementia (Alzheimer's disease on OpenNeuro \texttt{ds004504}) returns null ($d = -0.31$), with the diagnostic correctly predicting FAIL. CCR tracks \emph{functional} cortical-state dyscoherence on intact networks; it does not generalise to loss-of-integration through structural neurodegenerative atrophy. The pre-registered next experiment is direct application to a PTSD cohort with pre/post-therapy longitudinal EEG. This note is empirical-methods support for \emph{Life as Closure} \S 8 and \S 13.2, not a direct measurement of $P_{\Sigma_\mathcal{I}}$ on any substrate.
\end{abstract}
\paragraph{Programme map (review-packet 2026-05-22).}
This paper is part of the Existence / Life / Closure review packet. Stable packet references:
Paper I = \emph{Existence as Closure};
Paper II = \emph{Life as Closure};
Paper III = \emph{From Processing to Point of View};
Note A = \emph{Bioelectric Closure} (SL-001);
Note B = \emph{Cortical Phase Closure} (SL-002);
Note C = \emph{Closure-as-Distance} (cross-substrate methodology);
Note D = \emph{Trauma / Cortical Closure} (SL-005);
Note E = \emph{FlowIndex} (SL-006);
Note F = \emph{Dyadic Joint Meaning} (SL-007).
Extension IV = \emph{Cosmic Self-Pruning and Frame Recurrence} is withheld from the primary packet.

\paragraph{Status taxonomy.}
\textbf{Theorem / Definition} = formal structural claim under stated assumptions.
\textbf{Conditional Proposition} = bridge claim depending on H-RP, P-A, substrate mapping, or empirical interpretation.
\textbf{Empirical Result} = completed data analysis with stated effect size + significance.
\textbf{Empirical Proxy} = measurable first-order approximation, not the formal operator.
\textbf{Pre-registered Proposal} = future test, not evidence.



\section{Introduction}\label{sec:intro}

\emph{Life as Closure} \citep{SmartLifeAsClosure} \S 8 defines dyscoherence as $\mathcal{D}_\mathcal{I}(v) = \|v - P_{\Sigma_\mathcal{I}}(v)\|$, the distance between a living frame's substrate state $v$ and the projection onto its closure-stable self-model subspace $\Sigma_\mathcal{I}$. The companion paper identifies sustained elevated dyscoherence with the structural object that suffering instantiates under the Access Principle.

Operationalising this on real biological substrates requires a measurable, between-cohort inferential statistic. The companion methodology paper \citep{SmartClosureDistance} establishes closure-as-distance as a first-order empirical proxy for distance-from-constraint-manifold, conditional on the substrate satisfying the $D_1$--$D_5$ applicability diagnostic.

\bigskip
\noindent\fbox{\parbox{0.97\textwidth}{\textbf{Core claim.} The per-subject calibrated cortical closure residual (CCR) is a measurable first-order empirical proxy for chronic functional cortical dyscoherence on resting-state EEG. CCR is expected to discriminate clinical cohorts with elevated chronic dyscoherence (anxiety, depression, PTSD) from matched controls if and only if the substrate's between-cohort feature-shift satisfies the multi-channel integrative-organisation condition of the $D_1$--$D_5$ diagnostic. Outside that regime, CCR may detect a shift, but should not be expected to dominate single-axis baselines.}}
\bigskip

The contribution of this note is to:
\begin{enumerate}\itemsep0pt
\item Apply the SL-002 closure-residual pipeline verbatim, with no parameter retuning, between a single-substrate paired-design context (anaesthesia) and a between-cohort clinical-trait context (anxiety). The H4 default-weight robustness check is then meaningful.
\item Test CCR on real-data trait anxiety with the diagnostic prediction made before the empirical result is computed. (Confirmed PASS-predicted-PASS.)
\item Bound the methodology with a diagnostic-correctly-predicted FAIL on dementia, showing the proxy tracks functional, not structural, dyscoherence.
\item Pre-register a PTSD pre/post-therapy longitudinal test as the next, gold-standard replication.
\end{enumerate}

\section{Methods}\label{sec:methods}

\subsection{The cortical closure residual (CCR)}\label{sec:methods-ccr}

We use the SL-002 \texttt{compute\_features} pipeline \citep{SolutionLab002} unchanged. Briefly: source-localised or sensor-level EEG is decomposed into theta ($4$--$8$\,Hz), alpha ($8$--$13$\,Hz), beta ($15$--$30$\,Hz), and gamma ($30$--$45$\,Hz) bands via FIR bandpass filters (band edges match \citep{SolutionLab002} \S 2); per-band analytic signals are obtained by Hilbert transform; nine closure features are computed per $1$-second sliding window with $0.5$-s step: local phase smoothness, global Kuramoto order, global order standard deviation, task-related order, alpha and beta phase-gradient directionality (PGD), theta-gamma phase-amplitude coupling (PAC), alpha:beta cross-band PLV, and a surrogate $z$ statistic. The closure-residual vector for window $w$ in subject $s$ is the standardised concatenation of these nine features against the cohort's reference distribution. The L\textsuperscript{2}-norm $\|R_{s,w}\|$ is the per-window CCR; the per-subject mean is the inferential statistic compared across cohorts.

\paragraph{Reproducibility specifics (ds007609 anxiety analysis).}
\begin{itemize}\itemsep0pt
\item \textbf{Preprocessing.} Raw BrainVision data resampled to 250 Hz, $0.5$--$45$ Hz bandpass; bad-channel rejection by per-channel variance outlier ($> 3\sigma$ vs cohort mean); 30-s eyes-closed resting segment used.
\item \textbf{Channel subset.} The same $14$-electrode frontal subnetwork as in \citet{SolutionLab002} where present in the recording; full $32$-channel set for global-order features.
\item \textbf{Exclusion criteria.} Subjects with $<\!20$\,s of clean eyes-closed resting recording or $>\!25\%$ of windows flagged for artifact are excluded.
\item \textbf{Standardisation.} Per-subject reference distribution computed on the within-subject eyes-closed segment; features standardised against this within-subject reference, then per-feature standardised against the cohort's reference-condition mean and standard deviation (the same two-stage standardisation as SL-002).
\item \textbf{Quartile rule.} The high-STAI / low-STAI quartile contrast uses STAI trait scores in \texttt{phenotype/stai.tsv}; the upper quartile is defined as scores $\geq$ the cohort's $75$th-percentile, lower quartile as $\leq$ the $25$th-percentile.
\item \textbf{Frozen-script provenance.} The analysis script is \texttt{src/sl005\_ds007609\_adapter.py}; thresholds and feature weights were frozen before the analysis was run and have not been modified since.
\item \textbf{Diagnostic logs.} The full CAD-D1--D5-v1 diagnostic computation log is saved alongside the results in \texttt{results/sl005\_ds007609\_summary.json}.
\item \textbf{Statistical test.} Between-cohort effect size is Welch's Cohen's $d$ (un-pooled variance); $p$-value from Welch's $t$-test. Confidence intervals from non-parametric bootstrap with $10{,}000$ resamples are recorded in the results JSON.
\item \textbf{Hotelling $T^2$ implementation.} For the multivariate dispersion analysis, Hotelling $T^2$ is computed on the per-subject mean closure-residual vector using the cohort sample-covariance (Bessel-corrected, no shrinkage); per-feature contribution computed from the diagonalised quadratic form.
\end{itemize}

\paragraph{Standardisation reference.} For paired-design substrates (within-subject awake vs sed), the per-subject awake mean and standard deviation are used. For unpaired-design substrates (between-cohort, e.g.\ low-STAI vs high-STAI), the reference-cohort's per-feature mean and standard deviation are used.

\paragraph{Default weights.} All nine features carry equal weight in the L\textsuperscript{2}-norm. There is no per-feature weight calibration, no per-substrate retuning. The default-weight robustness check (H4) requires this.

\subsection{Applicability conditioning on $D_1$--$D_5$}\label{sec:methods-diagnostic}

CCR is conditional on the substrate satisfying the multi-channel integrative-organisation conditions of the $D_1$--$D_5$ diagnostic of \citep{SmartClosureDistance}. Before applying CCR to a real-data cohort, the $D_1$--$D_5$ diagnostic is computed on the per-subject signed-delta of the closure features. A substrate that fails the diagnostic is not expected to show a CCR shift even if functional dyscoherence is genuinely present at the structural level, because the L\textsuperscript{2}-norm-from-centroid proxy does not discriminate sub-threshold multi-channel shifts.

\subsection{Pre-registered hypotheses}\label{sec:methods-prereg}

Pre-registered hypotheses, frozen prior to real-data application:

\begin{itemize}\itemsep0pt
\item \textbf{H1}~(clinical-vs-healthy). Subjects with elevated chronic functional dyscoherence (clinical anxiety, depression, PTSD) show elevated CCR at rest relative to matched controls; threshold between-cohort Cohen's $d > 0.5$.
\item \textbf{H2}~(symptom gradient). Within-cohort, symptom severity (STAI trait score, BDI, PCL-5) correlates positively with CCR; threshold Spearman $\rho > 0.3$.
\item \textbf{H3}~(therapy reduction). Pre/post-therapy contrast in a clinical cohort shows reduced CCR post-therapy if therapy is effective; threshold paired $d_z > 0.5$. \emph{Pre-registered for future replication; not tested in the present paper}.
\item \textbf{H4}~(default-weight robustness across substrates). SL-002 \texttt{compute\_features} applied verbatim (no parameter retuning) returns a meaningful clinical effect; documented in code.
\item \textbf{H5}~(cross-cohort replication). The empirical result on the anchor cohort replicates on an independent cohort with a comparable clinical-trait scale; threshold $d > 0.5$ on second cohort, same default weights. \emph{Pre-registered for replication on ds004315 mood manipulation if data permits}.
\end{itemize}

\section{Results}\label{sec:results}

\subsection{Synthetic-data pre-registered validation}\label{sec:results-synthetic}

The synthetic PTSD-vs-healthy generator (in \texttt{src/sl005\_trauma\_closure.py}) constructs 100 subjects per cohort with simulated chronic-arousal dyscoherence injected as elevated global Kuramoto order, increased theta-gamma PAC, and reduced local phase smoothness in the PTSD cohort.

\begin{itemize}\itemsep0pt
\item H1 PTSD-vs-Healthy CCR elevation: $d = +4.31$ (threshold $> 0.5$) \textbf{PASS}.
\item H2 exposure-response: $d_z = +3.91$ (threshold $> 0.3$ as effect-size proxy) \textbf{PASS}.
\item H3 therapy reduction (within-subject): $d_z = +2.49$ (synthetic-only; threshold $> 0.5$) \textbf{PASS}.
\item H4 default-weight robustness: SL-002 \texttt{compute\_features} used verbatim → \textbf{PASS}.
\item H5 cross-cohort replication ($n=2$ generator seeds): $d = +41.38$ \textbf{PASS}.
\end{itemize}

Overall: $5/5$ pre-registered hypotheses pass on synthetic data. The pipeline is operational; the metric responds in the predicted direction to controlled dyscoherence perturbation. Synthetic-demo validation does not substitute for real-data validation.

\subsection{Real-data application: trait anxiety on \texttt{ds007609}}\label{sec:results-anxiety}

We applied CCR to OpenNeuro \texttt{ds007609} \citep{ShalamberidzeAnxiety2025}: resting-state $256$-channel HydroCel EEG ($n = 51$ subjects), restricted to a $32$-channel $10$--$10$ subset for direct comparability with the SL-002 empirical anchor. Subjects were grouped by STAI trait score: low-STAI (range $32$--$46$, $n = 12$) and high-STAI (range $57$--$69$, $n = 12$). Mean STAI trait of full cohort $= 51$.

\paragraph{H1 result.} Between-cohort calibrated CCR distance from the low-STAI centroid: Cohen's $d = +0.79$ (Welch $t = +1.93$, $p = 0.07$). Threshold $d > 0.5$ \textbf{PASS}. The strongest contributing features were global Kuramoto order ($d = +0.60$), theta-gamma PAC ($d = +0.47$), local phase smoothness ($d = +0.40$), and surrogate-z ($d = +0.40$).

\paragraph{Diagnostic prediction.} The $D_1$--$D_5$ diagnostic (unpaired-design extension) applied a priori to the cached feature matrix returned $D_1 = 0.43$ PASS, $D_4 = 3.68$ PASS, $D_5 = 0.00$ PASS, with $D_2 = 0$ and $D_3 = 0.64$ marginal under strict paired-design thresholds. The overall classification was PASS-leaning at the boundary; empirical $d = +0.79$ confirms PASS.

\paragraph{H2 within-cohort gradient.} Across the full $51$-subject cohort, Spearman correlation between STAI trait score and per-subject CCR is $\rho = +0.28$, below the pre-registered threshold of $\rho > 0.3$ but in the predicted direction. The threshold was set without prior knowledge of within-trait gradient strength; the within-cohort symptom-gradient correlation appears slightly weaker than the dichotomous quartile contrast, consistent with the noise floor of single-subject CCR estimates at $\sim 60$-second eyes-closed resting EEG.

\paragraph{H4 default-weight robustness.} SL-002 \texttt{compute\_features} was used verbatim with no parameter retuning. The H4 PASS at $d = +0.79$ between a within-subject anaesthesia paradigm (SL-002, ds005620) and a between-cohort trait-anxiety paradigm (this paper, ds007609) is a substantive cross-paradigm transfer of the closure-residual machinery without parameter calibration.

\subsection{The dispersion signature of chronic clinical states (two-statistic decomposition)}\label{sec:results-twostat}

A two-statistic decomposition (added in revision after observation of the SL-008 dementia null) clarifies what the CCR L\textsuperscript{2}-norm-from-centroid test detects on the anxiety cohort. Running both L\textsuperscript{2}-norm Cohen's $d$ (sensitive to cohort-mean shift and within-cohort dispersion expansion) and Hotelling's T\textsuperscript{2} (sensitive only to cohort-mean shift) on \texttt{ds007609} yields:

\begin{itemize}\itemsep0pt
\item L\textsuperscript{2}-norm Cohen's $d = +0.79$ (between-cohort, $p = 0.07$) -- \textbf{PASS}
\item Hotelling's T\textsuperscript{2} $p = 0.80$ -- \textbf{FAIL}
\item Per-feature variance ratio (high-STAI / low-STAI, geometric mean across features): $\mathbf{1.59}$ -- substantial dispersion increase
\end{itemize}

The empirical signal on this cohort is therefore \emph{not} a cohort-mean displacement (Hotelling's T\textsuperscript{2} is null); it is a \emph{subject-level dispersion expansion in cortical-state feature space}. High-STAI subjects do not all shift in the same direction; instead, each high-STAI subject occupies a different region of altered feature space, expanding the within-cohort variance. This is the precise multivariate signature consistent with \emph{Life as Closure}'s dyscoherence model \citep{SmartLifeAsClosure} \S 8: chronic dyscoherence increases each frame's distance from its closure-stable subspace $\Sigma_\mathcal{I}$, but different frames depart from $\Sigma_\mathcal{I}$ in different directions, producing dispersion expansion rather than coherent displacement. The L\textsuperscript{2}-norm-from-centroid test detects this; the cohort-mean test does not.

\paragraph{Implication.} The L\textsuperscript{2}-norm CCR test is the appropriate statistic for chronic functional dyscoherence claims because dyscoherence is heterogeneous across affected subjects (subject-specific direction of departure) while elevated in magnitude. Reporting Hotelling's T\textsuperscript{2} alongside identifies whether a substrate \emph{also} shows cohort-level convergent shift (consistent with cohort-wide pharmacological or structural change). The chronic-anxiety cohort shows dispersion-only signature; the dementia comparison (Section~\ref{sec:results-dementia}) shows the opposite signature (mean shift + dispersion decrease), which the framework correctly identifies as \emph{not} chronic functional dyscoherence.

\subsection{Independent-cohort replication: mood manipulation on \texttt{ds004315}}\label{sec:results-mood}

We applied CCR to OpenNeuro \texttt{ds004315} \citep{CavanaghMoodManipulation}: $51$ subjects randomly assigned to Sad-mood-induction ($n_{\mathrm{Sad}} = 25$) or Neutral-mood-induction ($n_{\mathrm{Neutral}} = 25$, with one excluded for missing data) prior to performing a Probabilistic Selection Task. Each subject was scored on the Beck Depression Inventory (BDI) at baseline (BDI range $0$--$11$ across the cohort, indicating predominantly non-clinical depression at baseline). The mood induction is randomised between subjects, so the unpaired-design diagnostic of \citep{SmartClosureDistance} applies. The closure feature pipeline of \citep{SolutionLab002} was applied verbatim with default weights, no parameter retuning between SL-005 anxiety (ds007609) and SL-005 mood manipulation (ds004315).

\paragraph{H1 result.} Between-group calibrated CCR distance from the Neutral cohort centroid: Cohen's $d = +0.70$ (Welch $t = +2.47$, $p = 0.018$). Threshold $d > 0.5$ \textbf{PASS}. The result is in the predicted direction (elevated CCR in the Sad-mood-induced cohort), with $p < 0.05$ on a between-group test at $n = 50$. Two-statistic decomposition: Hotelling's T\textsuperscript{2} $p = 0.93$ (\textbf{FAIL}); per-feature variance ratio (Sad / Neutral, geometric mean) $= \mathbf{1.61}$. Same signature as the anxiety cohort: dispersion increase, no cohort-mean shift. Acute sad-mood induction produces heterogeneous cortical-state expansion rather than convergent displacement, consistent with the chronic-dyscoherence interpretation of Section~\ref{sec:results-twostat}.

\paragraph{Diagnostic prediction.} The unpaired-design diagnostic returned $D_1 = 0.39$ PASS, $D_4 = 4.21$ PASS, $D_5 = 0.00$ PASS, with $D_2 = 0$ FAIL and $D_3 = 0.55$ FAIL under strict paired-design thresholds. Overall diagnostic prediction: FAIL. The empirical outcome is PASS at $d = +0.70$ with $p = 0.018$.

\paragraph{Diagnostic false negative.} This is the second observed case (after the trait-anxiety borderline result on \texttt{ds007609}) where the unpaired-design diagnostic predicts FAIL on a between-group cortical-EEG substrate that empirically passes the closure-as-distance test. The mechanistic explanation, consistent with the companion methodology paper's documented limitation, is that at small between-group cohort sizes ($n \leq 30$ per arm) the strict $\geq 80\%$ per-feature subject-sign-agreement threshold for $D_2$ becomes statistically unreachable even when a multi-feature shift is genuinely present; the empirical L\textsuperscript{2}-norm-from-centroid is more sensitive at this $n$ than the per-feature sign-test that drives $D_2$. The two CCR cohort replications (anxiety and mood manipulation) therefore both contribute to the per-design threshold-recalibration evidence base of \citep{SmartClosureDistance}.

\paragraph{H2 BDI gradient.} Across the full $50$-subject cohort, Spearman correlation between BDI score and per-subject CCR is $\rho = -0.11$ ($p = 0.46$). Below threshold and direction not aligned with prediction. The honest reading: BDI was a baseline-only screening scale; the experimental variable is randomised mood induction, not baseline depression. The H2 null is consistent with the mood induction (rather than baseline depression) being the carrier of the CCR effect on this cohort.

\paragraph{Per-feature breakdown.} Cohen's $d$ per closure feature (Sad vs Neutral): \texttt{pgd\_beta} $+0.27$, \texttt{pac\_theta\_gamma} $-0.22$, \texttt{task\_order} $-0.19$, \texttt{global\_order\_std} $-0.21$, \texttt{pgd\_alpha} $+0.18$, \texttt{global\_order} $+0.11$, others $|d| \leq 0.05$. Per-feature shifts are small individually; the L\textsuperscript{2}-norm aggregate is what discriminates --- consistent with the structural-multi-channel claim and the methodology paper's $\sqrt{k_\mathrm{eff}}$ prediction.

\subsection{Scope refinement: dementia as a diagnostic-correctly-predicted null}\label{sec:results-dementia}

To bound the methodology's scope, we also applied CCR to OpenNeuro \texttt{ds004504} \citep{MiltiadousDementia2023} ($n = 36$ Alzheimer's disease + $29$ healthy controls). The unpaired-design diagnostic returned $D_2 = 1$ FAIL, $D_3 = 0.66$ marginal: diagnostic predicts FAIL. Empirical CCR contrast: Cohen's $d = -0.31$ (Welch $t = -1.21$, $p = 0.23$) --- not significant. Diagnostic and empirical outcome agree.

The dementia null reveals a substantive scope refinement: CCR tracks \emph{functional} cortical-state dyscoherence on intact networks (anaesthesia, anxiety), not loss-of-integration through structural neurodegenerative atrophy. Substrates with progressive structural atrophy require methodology adapted to feature-asymmetric cancellation patterns, which is a separate research direction. The dementia null is treated as a scope-clarifying result, not a failure of CCR within its intended scope.

\subsection{Cross-substrate feature transfer}\label{sec:results-transfer}

The two-feature LOSO subset discovered on \texttt{ds005620} source-localised anaesthesia (beta-band PGD + alpha:beta cross-band PLV) was applied verbatim as a fixed feature subset to \texttt{ds007609} anxiety. Empirical Cohen's $d = +0.22$ on this fixed two-feature subset, compared to the in-cohort full-pipeline value $d = +0.79$ --- the methodology and discovery procedure transfer across substrates, but the discovered minimal feature subsets are substrate-specific. Anxiety's load-bearing features (global Kuramoto order, theta-gamma PAC, local smoothness) differ from anaesthesia's, consistent with the different molecular mechanisms involved (chronic sympathetic arousal modulation vs acute GABA/NMDA/K\textsuperscript{+}-channel pharmacology). CCR-as-methodology transfers; the substrate-specific signature is to be re-discovered per substrate.

\section{Relation to \emph{Life as Closure}}\label{sec:relation-life}

\emph{Life as Closure} \citep{SmartLifeAsClosure} \S 8 defines dyscoherence as $\mathcal{D}_\mathcal{I}(v) = \|v - P_{\Sigma_\mathcal{I}}(v)\|$. The present paper does not directly measure $P_{\Sigma_\mathcal{I}}$ on any substrate. It should therefore not be read as direct empirical validation of the full living-frame formalism. Instead, CCR supplies a measurable first-order empirical approximation: the within-cohort reference centroid acts as a cohort-aggregate proxy for $P_{\Sigma_\mathcal{I}}$; the L\textsuperscript{2}-distance to that centroid is the approximate dyscoherence.

The dichotomous high-vs-low trait-anxiety contrast on \texttt{ds007609}, with diagnostic-prediction PASS confirmed empirically at $d = +0.79$, instantiates the dyscoherence-as-chronic-clinical-state mapping of \emph{Life as Closure} \S 8 on real human resting-state EEG, with no parameter retuning across substrates. This is the empirical anchor for \emph{Life as Closure}'s dyscoherence claim under the restricted scope of \emph{functional} cortical-state integration.

\emph{Life as Closure} \S 13.2's claim that trauma persists structurally as elevated chronic dyscoherence and \S 13.3's claim that therapeutic intervention reduces it remain empirically untested in this paper. The pre-registered PTSD pre/post-therapy test (\S\ref{sec:external}) is the appropriate replication.

\section{Pre-registered external validation}\label{sec:external}

These tests are not presented as supporting evidence. They are pre-registered external replication targets designed to refine or falsify the metric.

\paragraph{PTSD vs healthy.} Required: $n \geq 14$ subjects per cohort, resting-state $\geq 32$-channel EEG, validated PTSD diagnostic (DSM-5 or PCL-5). Predict: $D_1$--$D_5$ passes (chronic functional dyscoherence at multi-feature level); $d > 0.7$ between-cohort. Falsifier: $D_1$--$D_5$ predicts PASS but $d < 0.3$ at $n \geq 14$.

\paragraph{PTSD pre/post-therapy longitudinal.} Required: $n \geq 14$ subjects with PTSD diagnosis, EEG recorded before and after evidence-based therapeutic intervention (CBT, EMDR, prolonged exposure). Predict: within-subject paired $d_z > 0.5$ in the direction of reduced CCR post-therapy. Falsifier: no reduction in CCR despite clinical recovery as measured by PCL-5.

\paragraph{Mood-manipulation between-group design (\texttt{ds004315}).} \emph{Completed in the present paper (\S\ref{sec:results-mood}); empirical $d = +0.70$, $p = 0.018$.} The result and methodology are reported above; this entry retains the original pre-registration text for reference.

\paragraph{Treatment-resistant depression vs first-line responders.} Required: depression EEG cohort stratified by treatment-response status. Predict: TRD subjects show elevated CCR relative to responders, consistent with sustained chronic dyscoherence.

\section{What this paper does not claim}\label{sec:nonclaims}

\begin{itemize}\itemsep0pt
\item Not that CCR is the closure operator or directly measures $P_{\Sigma_\mathcal{I}}$. It is a first-order empirical proxy for distance from a within-cohort reference centroid in feature space.
\item Not that elevated CCR proves a person is suffering. CCR is a between-cohort statistical signature, not an individual-subject diagnosis.
\item Not that the anxiety result generalises to all clinical conditions. The dementia null shows the methodology has a scope (functional vs structural).
\item Not that the synthetic-demo $5/5$ substitutes for clinical validation. The synthetic generator was constructed to inject chronic-arousal dyscoherence in directions the metric is designed to detect. The synthetic result demonstrates pipeline operability; it does not establish clinical sensitivity.
\item Not that the $D_1$--$D_5$ thresholds are universal constants. They are empirical applicability thresholds of version CAD-D1--D5-v1 of \citep{SmartClosureDistance}.
\item Not that this note proves the dyscoherence claim of \emph{Life as Closure}. It supplies a measurable boundary condition for testing that claim on real substrates.
\item Not that the pre-registered PTSD or mood-manipulation tests have been performed or are supported by data in this paper.
\end{itemize}

\section{Limitations and falsifiers}\label{sec:limits}

\begin{itemize}\itemsep0pt
\item \textbf{Cross-cohort replication is acute-mood, not trauma.} Cross-cohort replication is currently represented by the \texttt{ds004315} mood-manipulation cohort, which returned Cohen's $d = +0.70$, Welch $t = +2.47$, $p = 0.018$, but it is an acute mood-induction cohort rather than a clinical trauma cohort. The diagnostic predicted FAIL on this substrate under unpaired-design thresholds; the observed PASS is therefore a diagnostic false negative / calibration addendum (see \citet{SmartClosureDistance}), not a confirmed trauma replication. PTSD pre/post therapy with longitudinal EEG remains the gold-standard external validation.
\item \textbf{Anxiety as proxy for trauma.} STAI trait anxiety is a chronic-functional-arousal scale, not a trauma diagnostic. The mapping to \emph{Life as Closure} \S 13.2 trauma persistence is by proxy, not direct.
\item \textbf{Sample-size limitations on $D_3$ and within-cohort gradient.} At $n = 12$ per quartile cohort, within-cohort symptom-gradient correlation has limited power. The $\rho = +0.28$ is below threshold but consistent direction.
\item \textbf{H3 therapy-reduction is synthetic-only.} The therapy-reduction hypothesis has not been tested on any real-data substrate; pre-registered as a future external replication.
\item \textbf{CCR thresholds are empirical.} The pre-registered $d > 0.5$ and $\rho > 0.3$ thresholds are empirical applicability thresholds (metric version CCR-v1), not universal constants.
\item \textbf{Generalisation beyond cortical EEG anaesthesia + anxiety is open.} Additional clinical-state cohorts (depression, OCD, bipolar) require independent validation.
\end{itemize}

\section{Data and code availability}\label{sec:data}

Code is open-source under Apache-2.0; prose under CC BY 4.0. Implementations of the synthetic generator and real-data adapter are at \href{https://github.com/vfd-aria/solution-lab}{\texttt{github.com/vfd-aria/solution-lab}} under \texttt{experiments/005-vfd-trauma-cortical-closure}. The closure feature backend reuses \texttt{experiments/002-vfd-cortical-phase-closure} verbatim. Synthetic demo: \texttt{src/sl005\_trauma\_closure.py}. Real-data anxiety adapter: \texttt{src/sl005\_ds007609\_adapter.py}. Open data: \texttt{ds007609} \citep{ShalamberidzeAnxiety2025}, \texttt{ds004504} \citep{MiltiadousDementia2023} (for the dementia-null scope refinement). All pre-registered hypothesis thresholds were frozen prior to real-data application; the prediction-then-test order is preserved in the analysis script outputs.

\section{Conclusion}\label{sec:conclusion}

The contribution is not that CCR measures suffering in any individual subject. The contribution is the \emph{diagnostic boundary}: CCR is a measurable first-order empirical proxy for chronic functional dyscoherence on cortical EEG, conditional on the $D_1$--$D_5$ applicability diagnostic. Trait anxiety on \texttt{ds007609} returns $d = +0.79$ with diagnostic borderline-PASS-leaning; mood manipulation on \texttt{ds004315} returns $d = +0.70$ ($p = 0.018$) as an independent cross-cohort replication; dementia on \texttt{ds004504} returns null with diagnostic correctly predicting FAIL. The methodology tracks functional cortical-state dyscoherence on intact networks, not structural neurodegenerative loss. The PTSD pre/post-therapy longitudinal test is the appropriate next experiment. This note is empirical-methods support for \emph{Life as Closure}'s dyscoherence and trauma claims, not a philosophical proof of them.

\bibliographystyle{unsrtnat}
\bibliography{references}

\end{document}
